\startreportpage{25.4}
\templaterow{结\quad 论\\（结\quad 果）}{94}{
普通Ping共5对Echo请求/应答，Type/Code为8/0和0/0，Identifier为244、Sequence为0—4，56字节数据均原样返回。首对帧1、2的抓包RTT为5.429毫秒，Ping显示5.759毫秒；两者测量边界不同。IPv4首部校验和只覆盖IP首部，ICMP校验和覆盖整个ICMP消息，两者不可替代。

2500字节ICMP数据对应2508字节ICMP消息。请求与应答均在本机观察到两片：IPv4总长1500/1048，分片负载1480/1028，偏移原值0/185（字节偏移0/1480），MF为1/0。重组后ICMP数据仍为2500字节，首对双向数据一致。

TTL=1、2各5个请求均未捕获响应；TTL=3收到5个由\code{202.115.39.102}发送的Time Exceeded（11/0）；TTL=4首次收到目标Echo Reply。ICMP Traceroute共12个探测、TTL=1—4各3个，最终第4跳收到3个目标应答。首个差错响应为帧4，外层TTL64、引用原TTL1，匹配帧3的Sequence3。星号只能表示等待期内缺少匹配响应，不能据此断言后续不可达。
}
\templaterow{小\quad 结}{83}{
本次实验把命令参数、协议字段和原始字节相互核对。不同系统的Ping选项并不通用，必须先查本机帮助；ICMP数据长度、ICMP总长度、IPv4总长度和Ethernet帧长也应明确区分。Wireshark把时间戳从数据中单独解析后显示Data48，完整ICMP数据仍为56字节。

分片后每片各有IPv4首部，但只有第一片从ICMP首部开始；偏移185以8字节为单位，并非185字节。差错报文又含一层原IPv4和原ICMP首部，统计响应类型或读取TTL时要区分外层和引用层。路由器引用的TTL为1是本次实测事实，不应按理想图改成0。

本机抓包证实了当前目标、路径及时间窗口内的结果，不能定位某次无响应的具体丢失位置，也不能把初始TTL=4推广为其他主机或时刻的固定值。原始pcapng、命令日志、原始输出、逐帧字段和独立验证脚本随附件保留，报告可在Overleaf继续编辑并重新编译。
}
\TemplateReview
